% HD-S41. GPT-6 Astra Pro /root/sixfolds; inherited effort not exposed.
% 21 November 2026. All-rank actual normal-bireflection reduction.
\subsection{HD-S41: an irreducible normal bireflection subgroup in every rank}
\label{sec:irreducible-bireflection-all-ranks}

Let \(K\) be an imaginary quadratic field, let
\(G\subset GL_m(K)\), \(m>8\), be finite and absolutely
irreducible with exact character field \(K\), and use its
specified natural primitive rational group-algebra factor of a
connected full-\(G\) sphere cover. Retain the actual corner
\(K\)-action and the restricted theta polarization. Suppose
this factor is a signature-\((3,3)\) sixfold and its period
family has rank nine. Suppose also that \(G\) is reflection
free and has an absolutely irreducible normal subgroup
\(N\triangleleft G\) generated by actual bireflections.
Then, for \(K=\mathbb Q(i)\), its ordinary Hermitian
determinant is \([-1]\). For every imaginary quadratic
\(K\ne\mathbb Q(i),\mathbb Q(\sqrt{-3})\), no such
family occurs. There is no upper bound on \(m\) in these
assertions.

In particular the conclusion applies when the subgroup generated
by all actual bireflections is absolutely irreducible. The proof
does not infer that irreducibility from normality. Combining the
result with HD-S36 and HD-G27 also treats groups containing
reflections: for the Gaussian field only \([-1],[-3]\)
remain, and for the seven other class-number-one fields only the
previously proved split sources remain. The new reflection-free
assertion is stronger: none of those seven fields occurs under
the stated irreducible-normal-subgroup hypothesis.

\paragraph{1. Exact classification input and the three alternatives.}
The new higher-rank classification input is the bireflection specialization
of Guralnick--Saxl, \emph{Generation of finite almost simple
groups by conjugates}, J.~Algebra \textbf{268} (2003),
519--571, DOI \texttt{10.1016/S0021-8693(03)00182-0}:
Theorem~8.4, pp.~557--558, and Theorem~7.1,
pp.~548--549, with Lemma~7.4, p.~551.
The shorter Theorem~7.A, p.~521, explicitly retains scalar
multiples. These are the same published specializations read in
HD-S37 and HD-S40. The general numerical radical in indexed
text is not required here; every generator under consideration
has actual codimension two.

For the complex representation of \(N\), Theorem~8.4 gives
exactly three alternatives in degree greater than eight:
\(N\) is primitive with no invariant fixed tensor product;
it permutes \(m\) lines with permutation image \(A_m\)
or \(S_m\); or \(m=2r\) and it permutes \(r\)
two-dimensional blocks with permutation image \(S_r\).
The tensor-decomposable alternative in that theorem has degree
four and therefore does not occur here.

The last alternative has period rank at most seven by HD-G29,
now for every \(r\ge5\). Its characteristic kernel,
explicit central character projectors, individual Galois
descent, quotient curves, and primitive-factor isogenies are
all part of that proof. Thus it is excluded by the present
rank-nine assumption. The other two alternatives are treated
below, including every actual overgroup \(G\).

\paragraph{2. All the line projectors descend individually.}
In the line alternative let \(T\) be the kernel of the
permutation action of \(N\). The scalar-kernel argument of
HD-S40, paragraph2, applies in every degree \(m\ge7\):
an irreducible complex monomial group generated by actual
bireflections, with permutation image \(A_m\) or \(S_m\),
cannot have scalar diagonal kernel. Its proof uses the actual
weighted three-cycle maps and the perfect point stabilizer,
without a coefficient-field or bounded-degree hypothesis.
Thus \(T\) is abelian and nonscalar.

Every abelian normal subgroup of \(N\) has trivial image
in \(A_m\) or \(S_m\), since these groups have no
nontrivial abelian normal subgroup for \(m\ge5\).
Consequently it is contained in \(T\). Hence \(T\)
is the unique largest abelian normal subgroup of \(N\),
is characteristic in \(N\), and is normal in \(G\).
These are statements about the actual subgroups, not chosen
diagonal coordinates.

Let \(\chi_j:T\to\mathbb C^*\) be its characters on
the original lines. Equality of two characters is an invariant
equivalence relation on the labels. The natural \(A_m\)
or \(S_m\) action is two-transitive, so either all characters
are equal or they are pairwise distinct. The first would make
\(T\) scalar. Thus the characters are distinct and their
literal projectors are
\[
 p_j=\frac1{|T|}\sum_{t\in T}\chi_j(t)^{-1}\rho(t),
 \qquad p_jp_k=\delta_{jk}p_j,\qquad
 \sum_jp_j=1,\qquad\operatorname{rank}p_j=1.             \tag{S41.1}
\]
The identities follow directly by averaging the nontrivial
one-dimensional character \(\chi_k/\chi_j\).
For \(E=K(\zeta_{\exp T})\), all these projectors belong
to \(E\rho(T)\), whose full primitive central-idempotent
list they are. Every \(\gamma\in\operatorname{Gal}(E/K)\)
permutes them and commutes with conjugation by the original
\(K\)-matrices \(\rho(N)\). The permutation centralizer
of natural \(A_m\) or \(S_m\) is trivial: its point
stabilizer fixes only that point, so a commuting permutation
fixes every point. Therefore
\[
          p_j\in M_m(K),\qquad U_j=p_jK^m,
                    \dim_KU_j=1.                          \tag{S41.2}
\]
This is an explicit descent of each original line. Because
\(T\) is normal in \(G\), conjugation by every element
of \(G\) permutes the same projectors. Thus \(G\) has
an actual transitive \(K\)-line system.

For \(H=\operatorname{Stab}_G(U_1)\), a nonzero
\(v\in U_1\), and its actual character
\(\chi:H\to K^*\), the map
\[
 \operatorname{Ind}_H^G\chi\longrightarrow V_K,
       t_j\otimes1\longmapsto\rho(t_j)v                  \tag{S41.3}
\]
is an isomorphism on the coset bases. If \(K\) has only
the roots of unity \(\{\pm1\}\), this induced character
is rational-valued. That contradicts the exact original
character field \(K\). This excludes the line alternative
over every imaginary quadratic field except the Gaussian and
Eisenstein fields.

For \(K=\mathbb Q(i)\), the induced trace formula forces
\(\chi(H)=\mu_4\). HD-S34 now applies to exactly the
\(K\)-defined line system in (S41.3). Its actual cyclic
quotient is
\[
 Y=\widetilde C/H,\qquad
 X=\widetilde C/\ker\chi,\qquad X\longrightarrow Y.
                                                               \tag{S41.4}
\]
The scalar-normalized coefficient idempotent equals the lifted
\(\chi+\bar\chi\) idempotent after the
\(\ker\chi\)-average. Pullback to \(J\widetilde C\)
is an isogeny onto the selected line factor with multiplier
\(|\ker\chi|\), with rational inverse the norm divided by
\(|\ker\chi|\). The inverse action of an actual
\(h_0\in H\), \(\chi(h_0)=i\), is its Gaussian corner
action; all these maps intertwine that action. The parabolic
cup computation of HD-S34 gives ordinary determinant \([-1]\).
The sixth power of every displayed rational polarization
multiplier is a Gaussian norm. Paragraph6 below treats the
originally specified primitive line as well.

\paragraph{3. The primitive alternative and its exceptional degree.}
If \(N\) is complex primitive, then \(G\) is primitive:
any \(G\)-block system would also be an \(N\)-block
system, and irreducibility makes \(N\) transitive on its
blocks. An invariant fixed tensor product for \(G\) would
also be one for \(N\). Therefore the hypotheses of
Theorem~7.1 apply to \(G\). In characteristic zero and
degree greater than eight, the theorem gives the projectively
deleted permutation alternative, or the degree-ten normalizer
of the specified irreducible \(U_5(2)\) module.

The latter is excluded over every imaginary quadratic field.
The actual degree-ten representation has a unique invariant
nondegenerate alternating form, verified with the primary
ATLAS model in HD-S37, paragraph4. Every normalizer preserves
its line, with multiplier \(\lambda(g)\). For the actual
inertia tuple, irreducibility and the sixfold multiplicity give
\[
       \sum_{j=1}^b\operatorname{rank}(g_j-1)=26,
                   b\ge12,\qquad c_j\ge2.                 \tag{S41.5}
\]
Thus \(c_j\le4\). A nontrivial multiplier makes the fixed
space isotropic for the alternating form. Its dimension would
be at least six in a ten-dimensional symplectic space, which
is impossible. Every inertia multiplier is therefore one;
the inertia generate \(G\), so its representation is
self-dual. Finite-order dual eigenvalues are inverse roots of
unity, so every character value is real. A real value in an
imaginary quadratic field is rational, contradicting its
specified exact character field. This checks all normalizers
in the exceptional alternative, without replacing them by
scalar extensions.

\paragraph{4. A direct proof of the deleted-permutation normalizer.}
Put \(s=m+1\ge10\) and
\[
       V_0=\{x\in\mathbb C^s:\sum_i x_i=0\},\qquad
       Q(x)=\sum_i x_i^2,\qquad F(x)=\sum_i x_i^3.          \tag{S41.6}
\]
The complete complex linear normalizer of the standard
\(A_s\) on \(V_0\) is
\[
          N_{GL(V_0)}(A_s)=\mathbb C^*\,S_s.               \tag{S41.7}
\]
Here is a proof identifying the actual permutation map.
In degrees at most three, \(A_s\) and \(S_s\) have the
same monomial orbits on \(\mathbb C^s\): a monomial uses
at most three labels, so an odd permutation can be corrected
by a transposition on two unused labels. An invariant form on
\(V_0\) extends invariantly by the rational projection
\(x\mapsto x-s^{-1}(\sum_i x_i)\boldsymbol1\).
It follows from the degree-two and degree-three symmetric
monomial bases that the invariant quadratic and cubic spaces
on \(V_0\) are precisely \(\mathbb C Q\) and
\(\mathbb C F\), respectively. Both forms are nonzero,
and \(Q\) is nondegenerate.

A linear normalizer \(A\) therefore preserves these two
lines. The Hessian of \(F\) at \(v\in V_0\) is
\[
                 H_v(w,z)=6\sum_i v_iw_iz_i,
                                \qquad w,z\in V_0.         \tag{S41.8}
\]
Its rank is preserved under the actual transformation \(v\mapsto Av\).
Write \(k\) for the number of nonzero entries of \(v\).
For \(k<s\), a kernel vector satisfies
\(v_iw_i=t\) for every \(i\); a zero entry forces
\(t=0\). The kernel is the vectors on the \(s-k\)
zero coordinates with sum zero, of dimension \(s-k-1\).
Thus \(\operatorname{rank}H_v=k\). For \(k=s\), the
same equations give \(w_i=t/v_i\); the kernel has dimension
zero or one according as \(\sum_i v_i^{-1}\) is nonzero
or zero. Its rank is then \(s-1\) or \(s-2\), both
greater than two. Consequently the nonzero rank-two Hessian
locus is exactly
\[
                 \bigcup_{i<j}\mathbb C^*(e_i-e_j).       \tag{S41.9}
\]
The normalizer permutes these root lines and preserves their
\(Q\)-orthogonality. Distinct lines are orthogonal precisely
when their two-element supports are disjoint. Thus its induced
permutation is an automorphism of the edge-intersection graph
of the complete graph on \(s\) vertices.

The maximal pairwise-intersecting edge sets of size at least
four are exactly the full stars. Indeed, after two edges
\(ij,ik\), an edge not containing \(i\) must be \(jk\);
these three edges admit no fourth edge meeting all three.
This identifies stars intrinsically and gives a permutation
\(\pi\in S_s\) of the vertices. The intersection of the
stars at \(i,j\) is their single edge, so the root-line
permutation is precisely that of \(P_\pi\).
Now \(P_\pi^{-1}A\) fixes every root line. On the basis
\(r_i=e_i-e_s\), write its action as \(r_i\mapsto a_i r_i\).
The line of \(r_i-r_j\) is also fixed, forcing
\(a_i=a_j\) for every pair. The map is scalar. This proves
(S41.7) and the asserted explicit normalizer map, without an
automorphism-group classification of \(A_s\).

\paragraph{5. The actual rational alternating subgroup and the field map.}
In the projectively deleted alternative, the published module
description gives a complex intertwiner \(C_0\) with
\[
 C_0\rho(g)C_0^{-1}=a_gP_{\pi(g)},\qquad
            \pi(G)=A_s\ \hbox{or}\ S_s.                   \tag{S41.10}
\]
This is precisely the scalar-extended module specification,
not an assertion that the whole original group already has
those rational matrices. The pair \((a_g,\pi(g))\) is
unique: a permutation acting as a scalar on the augmentation
module is the identity, by applying it to the root lines and
then to \(e_i-e_j\). Hence \(a\) and \(\pi\) are
actual homomorphisms in this complex coordinate model.

Commutators cancel all scalar phases. Their permutation images
generate \((A_s)'=(S_s)'=A_s\), and every scalar-free
commutator is the corresponding literal permutation matrix.
It follows that
\[
          C_0\rho(G')C_0^{-1}=\{P_\tau:\tau\in A_s\}.
                                                               \tag{S41.11}
\]
In particular \(G'\) is an actual characteristic subgroup,
with the specified rational standard character under this
isomorphism; there is no unspecified central covering group
left in (S41.11).

Fix the resulting isomorphism \(\psi:G'\to A_s\).
In the rational basis \(e_i-e_s\) of \(V_0\), solve
the finite \(K\)-linear equations
\[
                 C\rho(x)=P_{\psi(x)}C,
                                      \qquad x\in G'.       \tag{S41.12}
\]
The standard \(A_s\)-augmentation module is absolutely
irreducible: its permutation module has commutant dimension
two by two-transitivity, and splits as its constant line plus
the augmentation module. Thus the complex solution space of
(S41.12) has dimension one and contains the invertible
\(C_0\). Gaussian elimination over \(K\) has the same
rank as over \(\mathbb C\), so gives a nonzero \(K\)-solution
\(C\), which is a nonzero scalar multiple of \(C_0\) and
is invertible. This constructs an actual \(K\)-intertwiner.

Because \(G'\) is characteristic, each \(C\rho(g)C^{-1}\)
normalizes the rational standard \(A_s\). The direct
normalizer proof in paragraph4 gives
\[
               C\rho(g)C^{-1}=\lambda(g)P_{\pi(g)}.
                                                               \tag{S41.13}
\]
Multiplying by the rational inverse permutation matrix shows
\(\lambda(g)I\in M_m(K)\), so \(\lambda(g)\in K\).
The permutation representation on \(V_0\) has no nontrivial
scalar permutation, so these are literal homomorphisms. Finiteness
makes \(\lambda(g)\) a root of unity. Thus (S41.13) is an
actual faithful map
\[
                        G\longrightarrow\mu(K)\times S_s.
                                                               \tag{S41.14}
\]
For \(\mu(K)=\{\pm1\}\), all its traces are rational,
contradicting the exact field \(K\). For the Gaussian field,
(S41.14) is exactly the correlated scalar-permutation hypothesis
of HD-G26, in every coefficient rank. Independence of the two
projections is neither assumed nor needed. Its theorem gives
ordinary determinant \([-1]\) for a rank-nine period family.

\paragraph{6. The original corner and polarization are retained.}
For completeness, the representation changes and the convenient
primitive lines in the preceding applications correspond to
actual polarized isogenies. A \(K\)-intertwiner \(C\)
identifies the given invariant positive coefficient form with
\(a\) times the model form for \(a\in\mathbb Q_{>0}\),
by absolute irreducibility and comparison of their \(K\)-valued
Gram matrices. Keeping the vector \(Cw\) makes the normalized
matrix coefficient
\(h(w,\rho(g)w)/h(w,w)\) literally unchanged. Hence its
rational corner idempotent inside the same relabelled
\(\mathbb Q[G]\) is unchanged; it has the same actual
Jacobian image and corner field.

For a subsequent change from a nonzero \(K\)-vector \(w\)
to a convenient nonzero \(v\), put
\[
 \begin{gathered}
 q_w=h(w,w),\quad q_v=h(v,v),\quad
 T_{v,w}(x)=v\,h(w,x)/q_w,\\
 b_{v,w}=\frac{2m}{|G|}\sum_{g\in G}
 \operatorname{Re}\operatorname{tr}
       (T_{v,w}\rho(g^{-1}))\,g .
 \end{gathered}                                               \tag{S41.15}
\]
The coefficients are rational. Matrix-coefficient averaging
acts as \(T_{v,w}\) on \(\rho\), its conjugate on
\(\bar\rho\), and zero on the other rational blocks.
The two complex characters are distinct because the exact
character field is imaginary quadratic. Thus, with their
orthogonal primitive corners,
\[
     b_{v,w}=e_vb_{v,w}e_w,\qquad
     b_{v,w}^{\dagger}b_{v,w}=(q_v/q_w)e_w.                 \tag{S41.16}
\]
For an integer \(M>0\) clearing the coefficients,
\(a_{Mb_{v,w}}\) is an actual \(K\)-linear isogeny between
the selected abelian factors, with pullback polarization
multiplier \(M^2q_v/q_w\). Its rational inverse is obtained
from \(b_{v,w}^{\dagger}\) divided by \(q_v/q_w\),
with the corresponding integer denominator for \(M\).
The group action is covariant on Jacobians and its Rosati
adjoint is the inverse deck action, as in HD-G29; these are
the same literal group-algebra maps.

In the line case this carries the prescribed corner to the
one in (S41.4). In the deleted-permutation case it carries it
to the vector \(se_1-\boldsymbol1\) of HD-G26. That proof
uses the actual quotients by the point stabilizer, scalar
kernel, and their intersection; its complementary cyclic
factor is the image of its displayed self-adjoint projector,
and the quotient pullback is an isogeny onto the primitive
factor with the exact norm scalar. Those maps intertwine the
inverse cyclic action with the original corner field.
All the positive rational multipliers just displayed change
a rank-six Hermitian determinant by
\(t^6=N_{K/\mathbb Q}(t^3)\). They therefore preserve the
ordinary determinant class and, in families with fixed finite
labels and lattices, identify period ranks.

This proves all the assertions, using actual coefficient,
curve, corner and polarization morphisms. HD-S03a turns
ordinary class \([-1]\) in signature \((3,3)\) into the
explicit split rational Hermitian space. For the seven
class-number-one fields \(\mathbb Q(\sqrt{-d})\),
\(d=2,7,11,19,43,67,163\), the stronger reflection-free
nonexistence above requires no class-number calculation.
The remaining all-rank gate is the case where the subgroup
generated by the actual bireflections is trivial or is not
absolutely irreducible; this theorem does not replace that
case by the one just proved.

\paragraph{Author check and source scope.}
The complete proof has been checked through the explicit
projectors, the scalar-diagonal lemma's all-degree scope,
the Hessian kernel equations, root-line normalizer map,
actual commutator subgroup, \(K\)-linear intertwiner system,
and the polarized primitive-line operator. The source uses
only the original published bireflection specializations,
whose indexed full text was reopened. The separate numerical
threshold observation from Hall's original printed paper is
not a premise of HD-S41. No unproved Galois stability,
normal-closure irreducibility, or determinant inference from
dimension alone is used.
